Segundo o relatório Education at a Glance 2023, publicado pela \textit{Organisation for Economic Co-operation and Development}, responsável pela divulgação dos principais indicadores relacionados à educação, aproximadamente 25\% dos brasileiros ingressantes em um curso superior de bacharelado abandonam seus estudos no primeiro ano, apenas 38\% terminam a graduação dentro do prazo previsto, enquanto 51\% se mantêm sem diploma mesmo três anos após o prazo de graduação esperado, como evidenciado em [1]. Isso corrobora com o fato de que a evasão acadêmica constitui um dos problemas mais persistentes enfrentados dentro das instituições de ensino em todo o mundo, esses dados são observados não apenas no Brasil, mas também em diversos outros países.

 evasão no ensino superior é um problema complexo e multifacetado, impactando diretamente os índices institucionais e o desenvolvimento social. Diante desse cenário, a aplicação de técnicas de Inteligência Artificial tem se consolidado como uma ferramenta vital de gestão preditiva, permitindo a identificação precoce de discentes propensos a abandonar o curso a partir de suas características demográficas e de desempenho acadêmico [2]. Nesse contexto, o presente trabalho tem como objetivo aplicar as técnicas de Ciência de Dados com embasamento teórico fundamentado em [3] para classificar estudantes quanto ao risco de evasão ou ao sucesso acadêmico. A base de dados utilizada para este projeto é pública, disponibilizada online pelo repositório \textit{University of California, Irvine} (UCI) [4], e originalmente feita como um estudo de caso de um Instituto Politécnico de Portalegre (IPP), em Portugal [5].
